% 第15批：《大学物理学学习指导》第5版自测题3“相对论”
% 选择3.1.1—3.1.10，原书印刷54—56页／PDF物理62—64页。
\subsection{自测题3：速度合成、时钟与相对论能量选择题}

\sourcetest{3.1.1}{54}{62}
\problemtext 长度为$L$的火筒匀速行驶，
人从其后端沿筒射出子弹，子弹\emph{相对火筒}的速率为$v_2$，
问火筒自有时钟测得子弹到前靶的时间。
\workings
筒身自身静止系中两端相距固有长度$L$，
子弹相对筒速率已知$v_2$，
\[
\boxed{\Delta t_{\rm tube}=\frac L{v_2}}.
\]
选\textbf{B}。
\crosscheck 不应将火筒对地速率$v_1$与子弹对筒速率
简单相加用于\emph{筒自身时钟}的传播时间。

\sourcetest{3.1.2}{54}{62}
\problemtext 飞船在地面系以$v$飞行，船首向船尾发送光信号，
飞船上接收者的本地时钟记录两事件相隔$\Delta t$，
求飞船固有长度。
\workings
船自身系两端静止，真空光速$c$，故
\[
\boxed{L_0=c\Delta t}.
\]
选\textbf{A}。
\crosscheck 信号发射在船首、接收在船尾，
两事件不是同地，故$\Delta t$不是这对事件的严格固有时，
虽然船的\emph{长度}是固有长度。

\sourcetest{3.1.3}{54}{62}
\problemtext 直尺静止于$K'$系，与$x'$轴成45°；
$K'$以非零速率沿$x$正向相对$K$运动，
求$K$中尺的夹角相对45°是增大还是减小。
\workings
同步测长时，沿运动方向长度缩短$\Delta x=\Delta x'/\gamma$，
垂直方向不变$\Delta y=\Delta y'$，
\[
\tan\theta=\frac{\Delta y}{\Delta x}
=\gamma\tan45^\circ=\gamma>1.
\]
故$\boxed{\theta>45^\circ}$，选\textbf{A}。
\crosscheck 夹角改变与运动方向的\emph{正负号}无关，只与$\gamma$有关。

\sourcetest{3.1.4}{54--55}{62--63}
\problemtext 判断：一惯性系中同地同时发生的两事件，
及另一对同一时刻却不同地的两事件，
在另一运动惯性系中是否仍可同时。
\workings
同地同时就是同一时空事件，
其$\Delta t=\Delta x=0$在任意惯性变换后仍为0。
对\emph{不同地}但同一时刻的事件，
沿分离方向作一维洛伦兹变换
$\Delta t'=\gamma(0-v\Delta x/c^2)\ne0$
（$v\Delta x\ne0$），
故通常失去同时性。按原题默认的一维相对运动约定，
选\textbf{A（第一组同时、第二组不同）}。
\crosscheck 若两个事件的三维空间位移
严格垂直于参考系相对速度，
仍有$\bm v\cdot\Delta\bm r=0$而继续同时。
所以该选择结论需要\emph{沿事件分离方向有非零速度分量}的约定。

\sourcetest{3.1.5}{55}{63}
\problemtext 电子静能约$0.5$MeV，动能约$0.25$MeV，
求其速率与光速的比。
\workings
$(\gamma-1)m_ec^2=0.25$MeV，$m_ec^2=0.50$MeV，
故$\gamma=1.5$，
\[
\boxed{\beta=\frac vc=\sqrt{1-\frac1{1.5^2}}
=\frac{\sqrt5}{3}\simeq0.745}.
\]
选\textbf{C（约$0.75c$）}。
\crosscheck 经典式$K=\tfrac12mv^2$在如此高能下不可使用。

\sourcetest{3.1.6}{55}{63}
\problemtext 关于狭义相对论的四条判断：
有静质量粒子速度限制、测量时间长度依参考系变化、
同时性的绝对性、相对运动时钟的时间延缓。
\workings
有静质量物体不能加速达到或超过$c$，命题(1)真；
长度和时间的比较可能依惯性系改变，(2)真；
分离两事件的同时性不是绝对的，(3)假；
相对运动钟的坐标时间间隔相对固有时可延长，(4)真。
选\textbf{B（1、2、4）}。
\crosscheck 原书术语“质量随运动状态改变”应在现代写法中
区分静质量不变与总能量$E=\gamma m_0c^2$可变。

\sourcetest{3.1.7}{55}{63}
\problemtext 飞船自身长度90m，相对地球以$0.8c$运动。
船尾发射光脉冲并于船头接收到，
求地面系中这两个事件的空间距离。
\workings
地面测得飞船长度
$L=L_0\sqrt{1-0.8^2}=90(0.6)=54$m。
在地面系，光追赶移动船头：
$c\,t=L+0.8c\,t$，
故$ct=54/(1-0.8)=270$m。
\[
\boxed{\Delta x_{\rm Earth}=270\,\mathrm m}.
\]
选\textbf{C}。
\crosscheck 题问的是\emph{发射事件到接收事件的空间间隔}，
并非地面同一时刻测量的飞船长度54m。

\sourcetest{3.1.8}{55}{63}
\problemtext 电子以速率$0.99c$运动，静能取0.51MeV，
求其相对论动能。
\workings
\[
\boxed{K=(\gamma-1)m_ec^2
=\left(\frac1{\sqrt{1-0.99^2}}-1\right)(0.51)
\simeq3.10\,\mathrm{MeV}}.
\]
选\textbf{C}。
\crosscheck 粒子速度接近光速时，经典动能公式严重低估。

\sourcetest{3.1.9}{55}{63}
\problemtext 年发电量为100亿千瓦时的核电站，
若电能完全来自核材料静质量减少，
估算所对应的质量亏损。
\workings
100亿kWh=$10^{10}$kWh，换算能量
$E=3.6\times10^{16}$J。
按$E=\Delta mc^2$，
\[
\boxed{\Delta m=\frac{3.6\times10^{16}}
{(3.0\times10^8)^2}=0.40\,\mathrm{kg}}.
\]
选\textbf{A}。
\crosscheck 这是假设100%质能转化为电能的理想质量差，
并非实际消耗的燃料总质量。

\sourcetest{3.1.10}{55--56}{63--64}
\problemtext 两个静质量均为$m_0$的粒子在实验室系
沿一条直线以速度$\pm v$相向运动，
完全非弹性碰合为一个静止复合粒子，
求它的静质量$M_0$。
\workings
初始总动量零，初总能量为$2\gamma m_0c^2$；
碰合后的复合物静止，故
$M_0c^2=2\gamma m_0c^2$，
\[
\boxed{M_0=\frac{2m_0}{\sqrt{1-v^2/c^2}}}.
\]
选\textbf{D}。
\crosscheck 与\emph{同速同向}合并仅得$2m_0$不同；
这里初始的相对动能转化成复合粒子的内能，
增加系统静质量。
